\clearpage
\VerdierPhysicalPageAnchor{0023}
$\mathcal A$ は十分な入射対象をもつ。したがって全導来関手は、とりわけ
「$F$ の右超導来関手」を復元できるが、それ以上の情報を与える。たとえば、
アーベル圏の間の加法的関手 $G:\mathcal B\to\mathcal C$ があれば、全導来
関手を合成できる：

\VerdierSemanticAnchor{intro-3-derived-composition}
\[
\mathrm R G\mathbin{\cdot}\mathrm R F:
\mathbf D^*(\mathcal A)\longrightarrow\mathbf D(\mathcal C).
\]
通常の場合、関手 $\mathrm R G\mathbin{\cdot}\mathrm R F$ は $GF$ の
全導来関手と標準的に同型である：

\VerdierSemanticAnchor{intro-3-derived-composition-isomorphism}
\[
\mathrm R(GF)\xrightarrow{\sim}
\mathrm R G\mathbin{\cdot}\mathrm R F,
\]
この合成公式は、われわれの言葉で\emph{合成関手のスペクトル系列}を表す。
しかもこれはより精密である。というのも、$\mathrm R GF(X)$
（$X$ は $\mathbf D^*(\mathcal A)$ の対象）のコホモロジー対象について
若干の情報を与えるだけでなく、それらを生み出す複体を「擬同型を除いて」
完全に記述するからである。この複体が分かれば、たとえばさらに別の全導来
関手と合成できる。

\VerdierSemanticAnchor{intro-4}
\noindent\textbf{4.}\quad
一般に、この形式体系では「擬同型を除いた」複体そのものを直接扱うことに
なり、それをコホモロジー対象だけで置き換えないよう注意する。（計算や
議論の末に求める「不変量」は多くの場合、複体のコホモロジー対象であるが、
常にそうとは限らない。スペクトル系列の微分も興味深い「不変量」である。）

\VerdierSemanticAnchor{intro-4-derived-relations}
擬同型を除いた複体の水準、全導来関手の水準、あるいはより一般に導来圏の
間の関手の水準で、さまざまな関係（同型関係、双対型定理のための随伴関係
など）を確立する。純粋に自動的な手続き（スペクトル系列）により、これら
の関係は複体のコホモロジー対象、それらが備える自然なフィルトレーション
などについて情報を与える。

\VerdierSemanticAnchor{intro-4-technical-intermediaries}
伝統的な導来関手の値、すなわちコホモロジー対象だけに関心がある場合で
さえ、導来圏における関係は、これらの関手を研究し重要な関係を定式化する
ための不可欠な技術的媒介となることが多い。

\VerdierSemanticAnchor{intro-4-simpler-relations}
最後に、本質的な利点として、導来圏で確立される関係は、そこからコホモロ
ジー対象について導かれる関係よりも、見かけ上も実際にも簡潔である。
